\chapter{The Thesis}
\label{c:thesis}

\begin{glance}
Between roughly 2005 and 2025 China fought and won three complete technology wars---solar photovoltaics, lithium-ion batteries, electric vehicles---by executing a single, learnable playbook, and artificial intelligence is its fourth run. The mechanism is not to match the West chokepoint for chokepoint but to engineer the chokepoints out of the system: open-weight frontier models at marginal prices, an open hardware ecosystem on an unrevocable instruction set, a memory-capacity arbitrage substituting commodity DRAM for scarce HBM, and a domestic semiconductor base. The connective tissue is openness deployed as a weapon against the layers where Western incumbents collect their rents. The argument is decomposed into four propositions, graded and falsified separately.
\end{glance}

\section{Three wars, one playbook}

In solar photovoltaics China went from assembler of imported cells to more than 90\% of every stage of the global supply chain, at module prices down more than 99\% from their 1976 level~\cite{owid-learning,pv-shares-2023}. In lithium-ion batteries---commercialized by Sony in 1991 and led for two decades by Japan and Korea---Chinese firms now hold roughly 69--75\% of global EV battery installations and over 80\% of cell manufacturing capacity~\cite{sne-2025,bnef-2025}. In electric vehicles China became at once the largest car market, the largest car exporter, and home to the largest EV maker~\cite{rhodium-evs,caam-2025}. Chapter~\ref{c:precedents} gives the campaigns in detail.

These were not three lucky outcomes but three executions of one playbook: patient state designation of a strategic industry; demand creation at home; subsidized scale-up far ahead of demand; brutal domestic competition as an evolutionary selection engine; capture of the entire upstream supply chain; a cost-down learning-curve war that bankrupts foreign incumbents; and finally the conversion of dominance into geopolitical leverage---export controls on the very technologies China was once accused of copying~\cite{mofcom-battery-controls,csis-rare-earth}.

\section{Why AI is the fourth war}

The prevailing Western assumption is that AI is different: that frontier AI depends on scarce, controllable inputs---leading-edge logic, high-bandwidth memory, proprietary interconnects, closed model weights---and that export controls on these chokepoints can durably preserve a lead. The book argues that it is failing on all fronts simultaneously, and failing in a characteristic way: not because China matches the West chokepoint for chokepoint, but because it is engineering the chokepoints out of the system.

\textbf{The model.} Chinese laboratories have converted the model layer into a commons: open weights on the top open positions of the independent indices, a majority of routed tokens, prices a small fraction of the closed frontier's, and a measured capability gap compressed to low single digits~\cite{hf-spring-2026,stanford-aiindex-2026}. Open weights at marginal cost do to the model layer what Chinese modules did to solar: they set a world price that rent-seeking business models cannot survive.

\textbf{The compute.} The state has designated the open RISC-V instruction set as national infrastructure, is standardizing matrix extensions and chip-to-chip interconnects across vendors, and has open-sourced the software stacks that would otherwise fragment the domestic ecosystem~\cite{riscv-policy,huawei-ub}. In parallel, a memory-capacity arbitrage---terabytes of cheap commodity DRAM for scarce HBM---attacks the assumption that frontier training requires centralized, HBM-bound superclusters at all.

\textbf{The semiconductor.} Beneath both, the fabs: 5\,nm-class logic without EUV, the world's fourth DRAM producer, three domestic toolmakers in the global top twenty~\cite{smic-5nm,cxmt-ipo}. The remaining gaps---EUV lithography, leading-edge HBM---are real, but the strategy does not require closing them on the West's terms; it requires making them irrelevant to the chosen architecture.

\section{Openness as a weapon}

The connective tissue is a deliberate inversion of the Western technology-rent model. Where the United States concentrated AI value in proprietary layers---CUDA, NVLink, closed frontier weights, x86/Arm licensing---China's strategy, formalized in the July 2025 Global AI Governance Action Plan and the ``AI+'' State Council initiative, is to make every one of them open, standardized and free~\cite{action-plan,ai-plus}. This is not altruism; it is the same move as pricing solar modules at \$0.09 per watt. An open layer generates no rent for the incumbent, but it generates \emph{adoption}, and adoption generates the scale, the ecosystem and the standards-setting power that constitute victory in the precedent industries. The two strategies are not symmetric opponents: one defends margins, the other destroys them.

\section{Four propositions, stated separately}
\label{c:propositions}

Rigor requires decomposing the thesis into four propositions that are related but logically independent---none proves the next, and each has its own evidence and its own failure modes.

\begin{description}
  \item[P1 --- Adoption.] China is winning open-weight model adoption. \emph{Confirmed by} download, derivative, and routed-token majorities persisting and extending into enterprise spend. \emph{Falsified by} a durable reversal of the router and Hugging Face shares, or enterprise open-weight spend stagnating below about 15\% while usage plateaus. Current grade: strongly supported [V].
  \item[P2 --- Sovereign compute.] China is constructing a domestically sufficient training-and-serving stack. \emph{Confirmed by} a frontier-class model trained end-to-end on domestic silicon, packaging, and memory. \emph{Falsified by} persistent dependence on stockpiled or smuggled foreign accelerators and HBM past about 2028. Current grade: rapid progress, not yet achieved [P/A/M].
  \item[P3 --- Rent destruction.] Open-weight commoditization will compress the closed-frontier rents that finance Western laboratories. \emph{Confirmed by} closed-model revenue growth decelerating while capability parity holds. \emph{Falsified by} closed labs sustaining pricing power through service, distribution, and trust layers even at model parity. Current grade: mechanism operating, magnitude contested [V/S].
  \item[P4 --- Financial repricing.] The compression of P3, if it occurs, punctures the US AI capital structure. \emph{Confirmed by} the indicator table of Chapter~\ref{c:trade}. \emph{Falsified by} demand elasticity absorbing the price decline: usage growing faster than prices fall rescues the buildout even under P1--P3. Current grade: scenario judgment [S].
\end{description}

They are coupled but separable. P1 can hold while P2 fails (China dominating distribution on stockpiled foreign silicon); P3 can hold while P4 fails (margins compress while volume growth fills the data centers); and a US valuation correction could arrive from macro causes with none of P1--P3 mattering. The convergence argument is that the \emph{coupling} of P1$\times$P2---open weights meeting sovereign, capacity-first hardware---is what forces P3, with P4 as its priced-in consequence; Chapter~\ref{c:summary} states where each proposition stands at the currency date.

\section{Evidence grades}

Load-bearing claims are graded throughout, on the scale the front matter defines: \textbf{[V]} verified, \textbf{[P]} primary-source claim---reliable as to what was claimed, not what was achieved---\textbf{[A]} analyst estimate, \textbf{[M]} author-modeled, \textbf{[R]} roadmap, \textbf{[S]} scenario judgment. The full edition collects every headline claim in a graded ledger and a machine-readable claim register \fullapp{D, J}; Appendix~\ref{c:register} reproduces the forecasts and the quarterly dashboard.

\section{The theory, in brief}
\label{c:theory}

The full edition devotes a chapter to the machinery behind the analogy \full{2}; four ideas from it are needed to read what follows.

\textbf{A cost stack, not a price.} The cost of one completed useful task decomposes into amortized model creation, inference, integration and organizational adoption, complementary capital, and assurance. Only inference is falling at ten times per year; integration and adoption are not---which is why open-weight \emph{usage} share rose while open-weight \emph{spend} share fell---and assurance is a cost the commons has mostly not paid. Commoditizing inference therefore does not commoditize delivered work; it shifts the binding constraint to the terms that are not falling, so the rent-destruction claim is bounded to the model layer and remains a hypothesis at the service layer.

\textbf{The elasticity knife-edge.} With own-price elasticity $\varepsilon$, revenue moves as $\Delta\log R=(1-\varepsilon)\,\Delta\log P$: above one, price cuts pay for themselves; below one, falling prices shrink revenue; $\varepsilon=1$ is the boundary between the bull and bear cases for the entire Western AI complex. The honest position is that $\varepsilon$ is a portfolio, not a number---plausibly above one for coding agents, synthetic data and scientific inference, near or below one for consumer chat---and the leveraged Western structure needs it above one at the premium tier specifically. Graceful degradation is a real Chinese advantage; it is not a theorem of victory.

\textbf{Openness relocates rent; it does not abolish it.} A seller rationally prices at or below marginal cost whenever the extra volume raises revenue on a complement by more than the margin given up. Each side opens the layer where it is behind and expects rent in the layer it is best at: China opens weights, ISA, interconnect and framework and expects to recapture on hardware volume, domestic cloud, standards and dependence; the United States controls all four and expects premium capability, distribution, integration, trusted deployment and agent platforms. Two flywheels result, and each can steal the other's critical input---China steals volume, America steals value---and which theft dominates is the single most important unresolved empirical question in the field.

\textbf{Four layers of winning.} AI is a general-purpose technology, so a country can dominate its production and capture little of its value, because most value appears downstream. Production leadership is the only layer on which the United States is unambiguously ahead; the platform layer goes to whoever opens fastest; diffusion is where most realized value sits; productivity is the only layer that finally matters and the least observable in real time.

The thesis is therefore conditional, and Chapter~\ref{c:summary} states it in its final form with the evidence attached: whether Chinese cost reduction, adoption, standards formation and hardware sufficiency compound faster than US capability, distribution and monetization. This book argues that convergence is the way to bet.
